% !TeX root = ../../report.tex
\subsection{Artificial Neural Networks}\label{sec:ann}

Artificial neural networks (ANNs), and specifically feedforward multi-layer perceptrons (MLPs), are a standard baseline for regression on tabular biomedical data. They are trained via backpropagation to minimise a loss function such as mean squared error (MSE) for continuous targets like the natural log of the half-maximal inhibitory concentration (LN IC50) \citep{goodfellow2016,arrieta2020}. In drug-response prediction, a lower LN IC50 indicates that less drug is needed to inhibit cell growth, and therefore a more drug-sensitive cell line.

Despite their wide use, unregularised neural networks are prone to overfitting on high-dimensional genomic data, where a model can easily find correlations in the training data that do not generalise \citep{klauschen2024}. The risk is greatest when the number of input features outnumbers the number of response samples \citep{partin2023}. An overfitted network memorises training-specific noise rather than the underlying relationship between molecular profiles and drug response. This appears as a generalisation gap, a large difference between performance on training data and unseen data \citep{keskar2017}. A related failure is \emph{Clever Hans} behaviour, in which a network relies on shortcuts such as spurious correlations or dataset artefacts instead of genuine signal \citep{lapuschkin2019,anders2022}. Overfitting and shortcut learning are not the same, but this project treats a wide generalisation gap as a possible warning sign of fragile explanations.

Pahde et al.\ \citeyearpar{pahde2023} show that deep networks frequently learn spurious correlations from their training data, which poses particular risks in medical applications \citep{pahde2023}. Their Reveal to Revise (R2R) framework was developed for image classification. Explanation methods can also look convincing even when they do not reflect what a model has learned \citep{adebayo2018}, so a plausible explanation is not evidence of a sound model. In this project, an intentionally overfitted network is therefore used as a diagnostic tool: its generalisation gap is the measurable starting condition (the Fragility Baseline) against which the effect of SAM is tested. Whether a smaller gap yields better explanations is the hypothesis under test, not an established result.

\subsubsection{Deep Learning for Cancer Sensitivity Prediction}\label{sec:deep-learning-cancer-sensitivity}

The standard MLP baseline described above has a structural limitation on genomic data. It treats every input feature as an independent, interchangeable variable and ignores the biological organisation that governs drug response. Genes do not act in isolation; they operate within signalling pathways that collectively regulate cell proliferation, survival, and resistance to therapy. Several studies report improved generalisation when network or pathway knowledge is built into the model \citep{partin2023}, and DeepCCDS outperformed a range of deep learning and conventional machine-learning methods in its own benchmarks \citep{wu2025}.

DeepCCDS \citep{wu2025} is a framework of this kind. It predicts the drug sensitivity of cancer cell lines by characterising cancer-driver signals. Its design departs from the traditional approach of feeding raw gene expression directly into the network. Instead, the influence of known cancer-driver genes is propagated through a prior-knowledge protein-interaction network, and the resulting signal is summarised as the activity of 38 biologically meaningful KEGG signalling pathways \citep{wu2025}. Each cell line is therefore described by a compact, interpretable vector of pathway-activity scores. These scores are computed by single-sample gene set enrichment analysis (ssGSEA) as a preprocessing step and are supplied to the network directly; they are not learned by it.

In addition to pathway activity, DeepCCDS takes binary mutation profiles across 534 cancer-driver genes and 1024-bit Morgan chemical fingerprints, compressing each high-dimensional binary input via a dedicated self-supervised autoencoder into a 30-dimensional embedding \citep{wu2025}. The pathway vector (38 dimensions) and the two embeddings (30 dimensions each) are concatenated into a 98-dimensional representation. A feedforward regression head then maps this representation to a single continuous output, the natural log of the half-maximal inhibitory concentration (LN IC50). The whole network is trained end-to-end on a combined objective, in which the drug-response prediction error is minimised jointly with the two autoencoder reconstruction errors \citep{wu2025}.

DeepCCDS was developed and evaluated on 319,543 drug--cell-line pairs from GDSC. These comprise 866 cell lines and 413 drugs, and the study reports predictive performance superior to that of the state-of-the-art methods it was compared against \citep{wu2025}. The authors report that t-SNE visualisations of the learned representations separate drug-sensitive from drug-insensitive cell populations, which they interpret as evidence that the model captures biological structure rather than statistical noise. They also report clinical relevance: predicted LN IC50 values distinguished responsive from non-responsive patient--drug combinations in an independent cohort of solid-tumour samples from The Cancer Genome Atlas (TCGA) \citep{wu2025}.

DeepCCDS is well suited to this project for three reasons. First, its architecture is fully specified and its inputs are available, so the baseline can be reproduced faithfully and the same network can be trained under different optimisers (Adam and SAM). This allows a controlled comparison of optimiser effects, provided other training settings are held fixed (Section 5.1). Second, its pathway-level inputs make it partly interpretable by construction. The pathway that a rule or attribution refers to has direct biological meaning, in contrast to an individual neuron or anonymous expression value. This supports the biological verification stage of the pipeline. When symbolic rules are extracted from the model, the pathways they reference can be related to established cancer biology, and pathway features carry no per-gene expression values, so the enrichment test was run on the mutation genes referenced by the rules (Section 6.2.5). Third, its compact, human-readable pathway representation supports the attribution audit, in which the model's reliance on individual pathways can be inspected and reported in biological terms.

It should be noted that the interpretability of DeepCCDS is a property of its inputs and not a guarantee about what the trained network has learned. The fact that pathway features are meaningful does not establish that the model uses them for the right reasons. This distinction is the point at which the Red XAI approach applies: a model with biologically meaningful inputs can still overfit and rely on unreliable shortcuts, so its explanations must be audited before they are trusted.

\subsubsection{Overfitting in Deep Models on Genomic Data and the Clever Hans Problem}\label{sec:overfitting-clever-hans}

A model exhibits Clever Hans behaviour when it achieves high performance by exploiting spurious correlations or dataset-specific artefacts rather than the underlying signal \citep{lapuschkin2019,anders2022}. Pharmacogenomic data are prone to such conditions. Partin et al.\ \citeyearpar{partin2023} note that, given the size of drug-response datasets, the number of input features can outnumber the number of response samples, which can lead to overfitting \citep{partin2023}. In this setting, unregularised networks tend to fit training-specific noise, settle in sharp regions of the loss landscape \citep{li2018}, and generalise poorly \citep{keskar2017}. The result is a generalisation gap: the difference between performance on training data and held-out data (Section 5.1.2). Overfitting and Clever Hans behaviour are distinct. Standard performance metrics can be oblivious to shortcut behaviour \citep{lapuschkin2019}, and because systematic spurious correlations may be present across the full dataset, including data reserved for testing, Clever Hans effects cannot be ruled out by test performance alone.

Reducing input dimensionality, as DeepCCDS does, does not remove this risk. The 319,543 drug--cell-line pairs in GDSC are not independent observations. Each cell line (866) is paired with many drugs, and each drug (413) with many cell lines, so the number of biologically independent units is much closer to 866 and 413 than to the number of pairs. The training, validation, and test sets in this project were drawn at random by pair, so cell lines and drugs appear in all three. The measured gap therefore reflects generalisation to unseen pairs, not unseen cell lines or drugs, and is likely to understate model fragility.

The most relevant methodological precedent is Reveal to Revise (R2R) \citep{pahde2023}, an iterative cycle of identifying spurious model behaviour, localising responsible artefacts, correcting the model, and re-evaluating it. Reliance on such artefacts is a particular risk in medical applications \citep{pahde2023,klauschen2024}. Adebayo et al.\ \citeyearpar{adebayo2018} add that explanation methods can produce convincing output that is largely independent of what the model has learned, so a plausible-looking explanation is not evidence that a model is sound \citep{adebayo2018}. R2R differs from this project in that it targets image data, uses attribution heatmaps, and corrects the model afterwards by fine-tuning. This project uses tabular genomic data and symbolic IF--THEN rules, and intervenes at the optimisation stage with SAM, aiming to prevent fragile explanations rather than repair them.

The baseline DeepCCDS model \citep{wu2025}, trained without explicit regularisation, serves as the Fragility Baseline. We hypothesise that the size of its generalisation gap indicates how fragile its explanations are (Section 1.4). Two checks establish the baseline: a quantitative measurement of the gap (Section 5.1.2) and a qualitative audit of the model's feature attributions in the spirit of Pahde et al.\ (Sections 5.1.2 and 6.2.4). The SAM+WD network is then compared against it to test whether a smaller gap improves the quality of the extracted rules.